\section{Introduction}
\label{sec:intro}

% Story lock (ARXIV.md). Author voice for the four-part stack is
% LOCKED in spirit --- polish only; do not delete or replace with the
% old ops.yaml claim list. Die figure lives in §impl (first cite).

% ---------- START CHECKING HERE ----------

The recent wave of Generative Artificial Intelligence ({Gen-AI}) models has
pushed large model operators toward designing custom accelerator
silicon~\cite{openai2026jalapeno,carter2026anthropic} in search of better performance per watt, 
a broader supplier base, and lower capital and operating cost.
That same demand has also drawn a crowded field of startups and open-source
efforts trying to fill the gaps at different levels of the stack.
Having worked inside several of them, I find the landscape sorts into
three camps:

\begin{itemize}
  \item \textbf{Pure-technology efforts}
        that invent a new way to do a hard kernel---an analog General
        Matrix Multiply ({GEMM}) or an in-memory compute fabric.
  \item \textbf{{RISC-V} intellectual-property ({IP}) efforts}
        that build high-end cores aimed at the datacenter.
        That market is still dominated by x86---about 87\% of
        datacenter server {CPU}s in early
        2025~\cite{mercury2025q1}---with Arm a distant second;
        a large share of the spend on this kind of effort goes to bringing those 
        cores up to run Linux, which is the wrong bar for bare-metal Artificial
        Intelligence ({AI}) acceleration anyway.
  \item \textbf{{AI} compiler efforts}
        that promise a holistic software stack, pushing yet another level of abstraction
        on hardware that is not yet shipping.
        Gemmini~\cite{genc2021gemmini} generates arrays;
        {LLVM}~\cite{lattner2004llvm}, {MLIR}~\cite{lattner2021mlir},
        and kin lower graphs onto backends someone else already built,
        with bloated and rapidly-changing codebases.
\end{itemize}

In general-purpose computing, the hardware-software\footnote{Unless otherwise specified, software means compiler} 
contract has always been an Instruction Set Architecture (ISA) and an Application Binary Interface (ABI), defining
what operations the CPU supports, general-purpose and purpose-specific registers, and how those registers should be
used by a compiler, and how memory accesses are treated.
The silent, underlying assumption was the homogeneity of the systems: every general-purpose computer built in the last 40+ years
had the same architectural shape. A CPU talking to a DRAM controller, with multiple level of caches interposed in between.
The functionality of the system could be extended by adding a card on an expansion slot, such as a Peripheral Component Interconnect Express ({PCIe}) network card or a Graphics Processing Unit ({GPU}).
The card, which we can think of as an accelerator per se, was its own microsystem that lived in a vacuum, interfacing with the main 
system by means of interrupts, PCIe transactions and kernel drivers.

This approach has served us well for the last five decades in computer organization and design, but it is now shortchanging us when we try to design a simple, and yet elusive accelerator system.
The heterogeneous nature of these systems makes it virtually impossible to have a clean, abstracted interface between hardware and software.
The same way Linux had to introduce Device Tree ({DT}) files to support the growing pool of heterogeneous, ARM-based Systems-on-Chip ({SoCs}), we need to rethink the shape
of the hardware--software interface in the era of AI accelerators.

That is why I built \tv{}: an open-source, {RISC-V}-based {AI}
accelerator infrastructure, composed of:
\begin{itemize}
  \item an extensible just-in-time ({JIT}) {AI} compiler (\pyv{});
  \item a runtime eXtensible {C} Library ({XCL}).
  \item a {RISC-V} core supporting {RV32IM} and the {RISC-V} Vector
        embedded integer profile ({Zve32x});
  \item an eXtensible Datapath Accelerator Interconnect ({XDAI});
\end{itemize}

\begin{figure}[t]
  \centering
  % Column-width PyVedas compile flow.
% Theme aligned with pyvedas_lower: white ground, black ink,
% gray only for emphasis; arrows on background layer.
\definecolor{tvink}{HTML}{111111}
\definecolor{tvmute}{HTML}{333333}
\definecolor{tvline}{HTML}{555555}
\definecolor{tvshade}{HTML}{EEEEEE}

\pgfdeclarelayer{backarrows}
\pgfsetlayers{backarrows,main}

\begin{tikzpicture}[
  >=Stealth,
  x=1mm, y=1mm,
  font=\sffamily\scriptsize,
  box/.style={
    draw=tvink,
    line width=0.6pt,
    rounded corners=1.5pt,
    fill=white,
    align=center,
    inner xsep=2.5pt,
    inner ysep=2.6pt,
    font=\sffamily\scriptsize,
  },
  hub/.style={
    box,
    line width=1.15pt,
    fill=tvshade,
    font=\sffamily\footnotesize\bfseries,
    inner xsep=4pt,
    inner ysep=3.0pt,
  },
  arr/.style={
    ->,
    draw=tvline,
    line width=0.55pt,
    shorten >=0.8pt,
    shorten <=0.8pt,
  },
]
\pgfmathsetmacro{\W}{0.92*\columnwidth/1mm}

% Thin frame only (no wash, no colored header bar)
\draw[tvink, line width=0.6pt, rounded corners=2.5pt]
  (0,-1.5) rectangle (\W,76.5);

\node[font=\sffamily\footnotesize\bfseries, text=tvink, anchor=west]
  at (2.0,73.8) {PyVedas compile flow};
\draw[tvink, line width=0.45pt]
  (2.0,71.6) -- (\W-2.0,71.6);

% --- row: inputs ---
\node[box, text width=0.38*\columnwidth] (model) at (0.26*\W, 60.5)
  {\textbf{PyTorch model}\\[0.35em]
   {\color{tvmute}graph}};

\node[box, text width=0.38*\columnwidth] (soc) at (0.74*\W, 60.5)
  {\textbf{SoC description}\\[0.35em]
   {\color{tvmute}XDAI + DAs}};

\node[box, font=\sffamily\scriptsize\bfseries, inner xsep=3pt, inner ysep=1.4pt]
  (wt) at (0.26*\W, 50.8) {weights};

% --- hub ---
\node[hub, minimum width=0.78*\columnwidth] (jit) at (0.5*\W, 40.2)
  {PyVedas\\[0.25em]
   {\normalfont\scriptsize\color{tvmute}extract · tile · lower}};

% --- C + link + ELF ---
\node[box, align=left, text width=0.72*\columnwidth,
  inner xsep=4pt, inner ysep=2.8pt] (cfile) at (0.5*\W, 27.6)
  {\textbf{Generated C}\\[0.2em]
   {\color{tvmute}calls mapped through the XCL}};

\node[box, minimum width=0.78*\columnwidth,
  inner xsep=4pt, inner ysep=2.8pt] (link) at (0.5*\W, 16.4)
  {\textbf{Link + stock RISC-V GCC}};

\node[hub, minimum width=0.78*\columnwidth,
  inner xsep=4pt, inner ysep=2.8pt] (elf) at (0.5*\W, 5.6)
  {Target ELF};

\begin{pgfonlayer}{backarrows}
  \draw[arr]
    (model.south) -- ++(0,-1.6) -| ([xshift=-3.5mm]jit.north);
  \draw[arr]
    (soc.south) -- ++(0,-1.6) -| ([xshift=3.5mm]jit.north);
  \draw[arr, densely dashed]
    (wt.south) -- ([xshift=-2.8mm]jit.north);
  \draw[arr]
    (jit) -- node[font=\sffamily\scriptsize, text=tvmute, right, inner sep=1.2pt]
      {C emit} (cfile);
  \draw[arr] (cfile) -- (link);
  \draw[arr] (link) -- (elf);
\end{pgfonlayer}

\end{tikzpicture}

  \caption{\pyv{} compile flow: model and {SoC} description in;
    graph and weights extracted; {C} against the {XCL}; stock
    {RISC-V} {GNU} toolchain out.}
  \label{fig:pyvedas-flow}
\end{figure}

The compiler, the runtime, the core, and that bus are the stack.

\section{Concept and Overview}

\pyv{} ingests a {PyTorch} model and a description of the target hardware, 
including any Datapath Accelerator ({DA}) present on the
{XDAI}.
It extracts the model's computational graph and its weights, then
lowers that graph into a {C} file whose calls target the {XCL}.
The {XCL} holds the mappings from {PyTorch} primitives to their
implementations.\footnote{For example, \texttt{aten.conv2d}.}
Those lowering passes extend to new {DA}s on the {XDAI}, including the
pre- and post-processing they need, such as tensor tiling and quantization
calibration.
The {C} is then linked against the {XCL} and built with the stock
{RISC-V} {GNU} toolchain.


To support a new {PyTorch} operator, \pyv{} needs a lowering candidate
and the {XCL} needs the matching {C} function---an \texttt{ops.yaml}
row, a codegen handler, and the sources it names.
To hang a new datapath, the same recipe plus the hardware {YAML} so
the {XDAI} decode and the {C} macros share one address window.
The {XCL} body then programs that window instead of unrolling the
work into {RISC-V}.

Each {DA} is a Memory-Mapped Input/Output ({MMIO}) device onto the core address space.
The reason why I choose to adopt an {MMIO} approach rather than implement
custom RISC-V instruction is mainly ease of integration, and to define a clear
boundary between Tiny-Vedas and third-party DAs. Supporting custom instructions
means maintaining a custom version of the {RISC-V} {GNU} Compiler Collection ({GCC}), or an out-of-tree version of
the RISC-V LLVM backend, which is simply impractical for most small teams.

As a running example throughout the paper, I choose an \texttt{int8} quantized version of 
{YOLOv3-Tiny}~\cite{redmon2018yolov3}, a popular computer vision model. The model is small enough to fit inside the
available Static Random-Access Memory ({SRAM}) on the Xilinx Alveo~{U280} Field-Programmable Gate Array ({FPGA}), but complex enough to present challenges when implementing
both software and hardware.

As a DA, I designed an \texttt{int8} $8{\times}8$ GEMM unit, 
representative of any third-party datapath accelerator you might want to attach to Tiny-Vedas,
and used it to offload the backbone of {YOLOv3-Tiny}.
